\chapter{Проектирование архитектуры программного комплекса}
\label{ch:design}

\section{Общая архитектура}

Система Fun-Walk построена по классической трёхуровневой схеме «клиент ---
прикладной сервер --- внешний сервис маршрутизации». На рис.~\ref{fig:architecture}
показано взаимодействие компонентов.

\begin{figure}[H]
\centering
\fbox{\parbox{0.92\textwidth}{\small\centering
\textbf{Frontend (React SPA, :5173)}\\
PreferencesPanel \textbar{} MapView \textbar{} RouteStats \textbar{} SavedRoutes\\
$\downarrow$ fetch('/api/...') $\downarrow$\\[0.4cm]
\textbf{Vite dev proxy} $\to$ \textbf{Backend (NestJS, :3000)}\\
RoutesController $\to$ RoutesService $\to$ RoutePlannerService\\
\hspace{2cm}$\to$ Dijkstra + POI scoring + OsrmRoutingService\\[0.4cm]
\textbf{OSRM} (router.project-osrm.org, profile foot)\\
\textbf{OpenStreetMap} tiles (Leaflet)
\vspace{0.6cm}}}
\caption{Архитектура системы Fun-Walk}
\label{fig:architecture}
\end{figure}

Клиент не содержит алгоритмической логики планирования --- все вычисления
выполняются на сервере, что упрощает сопровождение и позволяет в будущем заменить
алгоритм без обновления frontend.

\section{Модульная структура backend}

Backend организован в соответствии с принципами NestJS:

\begin{itemize}
  \item \texttt{AppModule} --- корневой модуль, подключает \texttt{ConfigModule}
        и \texttt{RoutesModule};
  \item \texttt{RoutesController} --- REST API, валидация DTO;
  \item \texttt{RoutesService} --- CRUD маршрутов, in-memory хранилище
        \texttt{Map<string, PlannedRoute>};
  \item \texttt{RoutePlannerService} --- оркестрация алгоритма (глава~\ref{ch:math});
  \item \texttt{OsrmRoutingService} --- HTTP-клиент OSRM.
\end{itemize}

Подмодуль \texttt{planner/} выделен явно для курсовой:
\begin{verbatim}
planner/
  route-planner.service.ts   # plan(), buildWeightedGraph(), calculateMetrics()
  math/geo.util.ts           # Haversine, projectOntoSegment, dotProduct
  math/poi-scoring.util.ts   # scorePois, selectCandidatePois
  graph/dijkstra.ts          # dijkstra()
  graph/graph.types.ts       # WeightedGraph, GraphNode, DijkstraResult
  routing/osrm-routing.service.ts
\end{verbatim}

\section{Структуры данных}

Ключевые типы определены в \texttt{route.interface.ts} и \texttt{graph.types.ts}.

\subsection{Точка и POI}

\begin{lstlisting}[language=TypeScript,caption={Фрагмент интерфейса PointOfInterest},label={lst:poi-interface}]
interface PointOfInterest {
  id: string;
  name: string;
  type: 'park' | 'green_zone' | 'bike_path'
      | 'waterfront' | 'square';
  location: LatLng;
  airQualityIndex: number;  // AQI
  noiseLevel: number;       // dB
  areaHa: number;
}
\end{lstlisting}

\subsection{Граф}

\begin{lstlisting}[language=TypeScript,caption={Типы взвешенного графа},label={lst:graph-types}]
interface GraphNode {
  id: string;
  location: LatLng;
  attractiveness: number;  // [0, 1]
  isPoi: boolean;
  poiName?: string;
}

interface WeightedGraph {
  nodes: Map<string, GraphNode>;
  adjacency: Map<string, { targetId: string; weight: number }[]>;
}
\end{lstlisting}

Использование \texttt{Map} обеспечивает $O(1)$ доступ к вершинам по идентификатору.

\subsection{Результат планирования}

Объект \texttt{PlannedRoute} включает поля для отчётности: \texttt{algorithm:
'dijkstra'}, \texttt{routingSource}, \texttt{graphStats: \{ nodes, edges, pathWeight \}}.

\section{REST API}

Контроллер \texttt{@Controller('api')} экспонирует endpoints (табл.~\ref{tab:api}).

\begin{table}[H]
\centering
\caption{REST API системы Fun-Walk}
\label{tab:api}
\begin{tabular}{@{}llp{7cm}@{}}
\toprule
\textbf{Метод} & \textbf{Путь} & \textbf{Назначение} \\
\midrule
GET & /api/health & Проверка работоспособности \\
GET & /api/defaults & Старт/финиш по умолчанию \\
GET & /api/poi & Список POI \\
GET & /api/routes & Сохранённые маршруты \\
GET & /api/routes/:id & Детали маршрута \\
POST & /api/routes/plan & Построение маршрута \\
DELETE & /api/routes/:id & Удаление \\
\bottomrule
\end{tabular}
\end{table}

DTO \texttt{PlanRouteDto} валидируется через \texttt{class-validator}:
$\varphi \in [-90, 90]$, $\lambda \in [-180, 180]$, preferences $\in [0, 10]$.

\section{Проектирование frontend}

Frontend --- одностраничное приложение (SPA) с компонентной декомпозицией:

\begin{itemize}
  \item \texttt{App.tsx} --- состояние приложения, вызовы API;
  \item \texttt{PreferencesPanel} --- шесть ползунков, режим выбора точки;
  \item \texttt{MapView} --- Leaflet, polyline, маркеры, fitBounds;
  \item \texttt{RouteStats} --- метрики и кольцевой индикатор score;
  \item \texttt{SavedRoutes} --- список маршрутов сессии.
\end{itemize}

Vite проксирует \texttt{/api/*} на \texttt{http://127.0.0.1:3000}, устраняя CORS
в режиме разработки. Production-сборка предполагает reverse proxy или настройку
\texttt{CORS\_ORIGINS} в \texttt{.env}.

\section{Конфигурация и развёртывание}

Переменные окружения (\texttt{backend/.env}):

\begin{table}[H]
\centering
\caption{Конфигурационные параметры}
\label{tab:env}
\begin{tabular}{@{}llp{6.5cm}@{}}
\toprule
\textbf{Переменная} & \textbf{По умолчанию} & \textbf{Описание} \\
\midrule
PORT & 3000 & Порт API \\
HOST & 0.0.0.0 & Адрес прослушивания \\
OSRM\_URL & router.project-osrm.org & Базовый URL OSRM \\
CORS\_ORIGINS & localhost:5173 & Разрешённые origin \\
\bottomrule
\end{tabular}
\end{table}

Корневой \texttt{package.json} использует \texttt{concurrently} для команды
\texttt{npm run dev}. Скрипты \texttt{scripts/dev.ps1}, \texttt{dev.sh}, \texttt{dev.bat}
автоматизируют установку и запуск на Windows и Linux.

\section{Диаграмма последовательности POST /api/routes/plan}

\begin{enumerate}
  \item Клиент отправляет JSON с \texttt{start}, \texttt{end}, \texttt{preferences};
  \item \texttt{RoutesController} валидирует DTO;
  \item \texttt{RoutesService.planRoute()} вызывает \texttt{RoutePlannerService.plan()};
  \item Planner: scoring $\to$ graph $\to$ Dijkstra $\to$ OSRM $\to$ metrics;
  \item Сервис сохраняет \texttt{PlannedRoute} в Map, возвращает UUID;
  \item Клиент отображает polyline и обновляет список маршрутов.
\end{enumerate}

\section{Ограничения проектных решений}

\begin{itemize}
  \item Хранение маршрутов in-memory --- данные теряются при перезапуске;
  \item POI статичны, не загружаются из Overpass API;
  \item AQI и шум --- заданные константы, не live-данные;
  \item Зависимость от публичного OSRM (rate limits, SLA).
\end{itemize}

Эти ограничения задокументированы как направления развития (глава~\ref{ch:conclusion}).

\section{Выводы по главе}

Спроектирована модульная клиент-серверная архитектура с чётким разделением
ответственности: UI, REST, алгоритмы, внешняя маршрутизация. Структуры данных
и API спроектированы с учётом требований курсовой отчётности (\texttt{graphStats},
\texttt{algorithm}).
